\documentclass[11pt,a4paper]{article}
\usepackage[utf8]{inputenc}
\usepackage[vietnamese]{babel}
\usepackage{amsmath,amssymb}
\usepackage[version=4]{mhchem}
\usepackage{geometry}
\geometry{top=1.5cm,bottom=1.2cm,left=1.4cm,right=1.4cm,headheight=28pt,headsep=12pt,footskip=18pt}
\usepackage{graphicx}
\usepackage{xcolor}
\usepackage{tabularx}
\usepackage{array}
\usepackage{tikz}
\usepackage{fancyhdr}
\usepackage{enumitem}
\usepackage{multido}
\usepackage{tcolorbox}
\usepackage{colortbl}

\definecolor{hfRed}{HTML}{D32F2F}
\definecolor{hfGreen}{HTML}{8BC390}
\definecolor{mainGreen}{HTML}{2E7D32}
\definecolor{solBlue}{HTML}{1565C0}
\definecolor{tableGreen}{HTML}{E8F5E9}

\usepackage{ex_test}

\pagestyle{fancy}
\fancyhf{}
\renewcommand{\headrulewidth}{0.75pt}
\renewcommand{\headrule}{\hbox to\headwidth{\color{hfGreen}\leaders\hrule height \headrulewidth\hfill}}
\renewcommand{\footrulewidth}{0.75pt}
\renewcommand{\footrule}{\hbox to\headwidth{\color{hfGreen}\leaders\hrule height \footrulewidth\hfill}}

\fancyhead[L]{\small\itshape\color{hfRed}Tài liệu lưu hành nội bộ}
\fancyhead[R]{\small\color{hfRed}Thầy \textbf{TRẦN VĂN THẠNH} - 0777.470.803}

\fancyfoot[L]{\footnotesize\itshape\color{hfRed}CS1: 6/15 Nguyễn Hoàng, P. Kim Long, TP Huế \quad|\quad CS2: 24 Đặng Thái Thân, TP Huế}
\fancyfoot[R]{\small\textbf{\color{hfRed}Trang \thepage}}

\setlength{\parindent}{0pt}
\setlength{\parskip}{2pt}

\begin{document}

\noindent
\begin{tabularx}{\textwidth}{@{}X r@{}}
\textbf{TRƯỜNG THPT LƯƠNG VINH} & \textbf{ĐỀ KIỂM TRA GIỮA HỌC KỲ I} \\
\textbf{ĐỀ CHÍNH THỨC} & \textbf{MÔN: HOÁ HỌC 12} \\
\textit{(Đề có 04 trang)} & \textbf{Thời gian làm bài: 45 PHÚT} \\
& \textbf{Mã đề 1204}
\end{tabularx}

\vspace{2pt}
\noindent\textbf{Họ và tên:} \dotfill \textbf{Lớp:} \dots\dots\dots\dots

\vspace{2pt}
\noindent\textit{* Cho biết nguyên tử khối của các nguyên tố: H = 1; C = 12; N = 14; O = 16.}
\vspace{2pt}
\hrule height 0.5pt
\vspace{6pt}

\noindent{\color{mainGreen}\textbf{\large PHẦN I. Câu hỏi trắc nghiệm nhiều phương án lựa chọn.}} \\
\textit{Thí sinh trả lời từ câu 1 đến câu 12. Mỗi câu hỏi thí sinh chỉ chọn một phương án.}
\vspace{4pt}
\begin{ex}
Glucose và fructose đều
\choice
{có nhóm \ce{-CH=O} trong phân tử.}
{thuộc loại disaccharide.}
{có phản ứng tráng bạc.}
{có công thức phân tử \ce{C6H10O5}.}
\loigiai{}
\end{ex}
\begin{ex}
Hợp chất \ce{C2H5NHC2H5} có tên là
\choice
{diethylamine.}
{dimethylamine.}
{ethylmethylamine.}
{propylamine.}
\loigiai{}
\end{ex}
\begin{ex}
Tinh bột thuộc loại polysaccharide, có nhiều trong các loại ngũ cốc như hạt lúa, hạt ngô, củ khoai, củ sắn, .... Công thức của tinh bột là
\choice
{\ce{C12H22O11}.}
{\ce{C2H4O2}.}
{\ce{C6H12O6}.}
{\ce{(C6H10O5)_n}.}
\loigiai{}
\end{ex}
\begin{ex}
Ester \ce{HCOOCH3} phản ứng với dung dịch \ce{NaOH} (đun nóng), sinh ra các sản phẩm hữu cơ là
\choice
{\ce{HCOOH} và \ce{CH3ONa}.}
{\ce{CH3ONa} và \ce{HCOONa}.}
{\ce{HCOONa} và \ce{CH3OH}.}
{\ce{CH3COOH} và \ce{CH3OH}.}
\loigiai{}
\end{ex}
\begin{ex}
Cho các chất: \ce{H2NCH2COOH} (1); \ce{C2H5COOH} (2); \ce{C2H5NH2} (3); \ce{H2NCH2CH2CH(NH2)COOH} (4); \ce{C6H5NH2} (5). Những chất vừa phản ứng được với acid vừa phản ứng được với base là
\choice
{(1), (2).}
{(2), (3).}
{(4), (5).}
{(1), (4).}
\loigiai{}
\end{ex}
\begin{ex}
Hợp chất nào sau đây được sử dụng làm chất giặt rửa tổng hợp?
\choice
{\ce{CH3[CH2]12COONa}.}
{\ce{CH3[CH2]11OSO3Na}.}
{\ce{CH3COONa}.}
{\ce{CH3[CH2]12COOCH3}.}
\loigiai{}
\end{ex}
\begin{ex}
Số nguyên tử hydrogen trong phân tử methyl formate là
\choice
{2.}
{6.}
{4.}
{8.}
\loigiai{}
\end{ex}
\begin{ex}
Phần trăm khối lượng nitrogen trong phân tử aniline bằng
\choice
{15,73\%.}
{15,05\%.}
{18,67\%.}
{12,96\%.}
\loigiai{}
\end{ex}
\begin{ex}
Dung dịch chất nào sau đây hòa tan \ce{Cu(OH)2}, thu được dung dịch có màu xanh lam?
\choice
{Propyl alcohol.}
{Propan-1,3-diol.}
{Albumin.}
{Fructose.}
\loigiai{}
\end{ex}
\begin{ex}
Một phân tử saccharose có
\choice
{một đơn vị $\beta$-glucose và một đơn vị $\alpha$-fructose.}
{hai đơn vị $\alpha$-glucose.}
{một đơn vị $\alpha$-glucose và một đơn vị $\beta$-fructose.}
{một đơn vị $\beta$-glucose và một đơn vị $\beta$-fructose.}
\loigiai{}
\end{ex}
\begin{ex}
Phát biểu nào sau đây đúng?
\choice
{Hai đơn vị glucose trong maltose liên kết với nhau bằng liên kết $\alpha-1,2-$glycoside.}
{Saccharose và maltose có cùng công thức phân tử.}
{Một phân tử saccharose gồm hai đơn vị $\alpha-$glucose.}
{Maltose không có nhóm -OH hemiacetal.}
\loigiai{}
\end{ex}
\begin{ex}
Phát biểu nào sau đây là \textbf{sai}?
\choice
{Dung dịch lysine làm xanh quỳ tím.}
{Aniline tác dụng với nước bromine tạo thành kết tủa trắng.}
{Dung dịch glycine không làm đổi màu quỳ tím.}
{Cho \ce{Cu(OH)2} vào dung dịch lòng trắng trứng thấy xuất hiện màu vàng.}
\loigiai{}
\end{ex}

\vspace{6pt}
\noindent{\color{mainGreen}\textbf{\large PHẦN II. Câu trắc nghiệm đúng sai.}} \\
\textit{Thí sinh trả lời từ câu 1 đến câu 2. Trong mỗi ý a), b), c), d) ở mỗi câu, thí sinh chọn đúng hoặc sai.}
\vspace{4pt}

\begin{ex}
Cho các carbohydrate: glucose, fructose, saccharose, maltose.
\choiceTF
\True Cả 4 chất ở điều kiện thường đều là chất rắn, vị ngọt, tan tốt trong nước.
\True Có 2 chất có nhóm -OH hemiacetal.
\True Có 2 cặp chất là đồng phân của nhau.
\True Có 3 chất vừa có cấu tạo dạng mạch hở, vừa có cấu tạo dạng mạch vòng.
\loigiai{}
\end{ex}

\begin{ex}
Cho công thức cấu tạo của peptide Y như sau:\\
\centerline{\ce{HOOC-CH2-CH2-CH(NH2)-CO-NH-CH2-CO-NH-CH(CH3)-COOH}}
\choiceTF
\True Tên gọi của Y là Glu-Gly-Ala.
1 mol Y phản ứng tối đa được với 3 mol \ce{NaOH}.
\True Thủy phân hoàn toàn Y thu được glutamic acid, glycine và alanine.
Phần trăm khối lượng của nitrogen (N) trong Y là 15,38\%.
\loigiai{}
\end{ex}

\vspace{6pt}
\noindent{\color{mainGreen}\textbf{\large PHẦN III. Câu trắc nghiệm yêu cầu trả lời ngắn.}} \\
\textit{Thí sinh trả lời từ câu 1 đến câu 4.}
\vspace{4pt}

\begin{ex}
Tiến hành thí nghiệm sau:\\
- Bước 1: Cho vào ống nghiệm 2 mL nước cất.\\
- Bước 2: Nhỏ tiếp vài giọt aniline vào ống nghiệm, sau đó nhúng giấy quỳ tím vào dung dịch trong ống nghiệm.\\
- Bước 3: Nhỏ tiếp 1 mL dung dịch \ce{HCl} đặc vào ống nghiệm.\\
Cho các phát biểu sau:\\
(1) Sau bước 2, chất lỏng bị phân thành 2 lớp, lớp dưới là aniline.\\
(2) Sau bước 2, giấy quỳ tím chuyển thành màu xanh.\\
(3) Sau bước 3, dung dịch thu được trong suốt.\\
(4) Sau bước 3, trong dung dịch có chứa muối phenylammonium chloride tan tốt trong nước.\\
(5) Ở bước 3, nếu thay \ce{HCl} bằng \ce{Br2} thì sẽ thấy xuất hiện kết tủa màu vàng.\\
Liệt kê các phát biểu \textbf{sai} theo thứ tự tăng dần (ví dụ 12, 13, ...).
\shortans{25}
\loigiai{
\textbf{Trả lời: 25}\\
- (2) Sai: Aniline có tính base rất yếu không đổi màu quỳ tím.\\
- (5) Sai: Tác dụng với bromine tạo kết tủa màu trắng, không phải màu vàng.\\
$\Rightarrow$ Phát biểu sai là (2) và (5) $\rightarrow$ 25.
}
\end{ex}

\begin{ex}
Cho các chất: methylamine, glycine, aniline, acetic acid, lysine, valine. Có bao nhiêu chất phản ứng với dung dịch \ce{NaOH} tạo muối?
\shortans{4}
\loigiai{
\textbf{Trả lời: 4}\\
Gồm 4 chất: glycine, acetic acid, lysine, valine (đều chứa nhóm carboxyl -COOH).
}
\end{ex}

\begin{ex}
Cho các chất: \ce{CH3[CH2]14COONa}, \ce{CH3[CH2]10CH2OSO3Na}, \ce{C2H5OSO3Na}, \ce{CH3[CH2]16COOK}, \ce{CH3[CH2]11C6H4SO3Na}. Có bao nhiêu chất là chất giặt rửa tổng hợp?
\shortans{2}
\loigiai{
\textbf{Trả lời: 2}\\
Gồm 2 chất: \ce{CH3[CH2]10CH2OSO3Na} (muối alkyl sulfate) và \ce{CH3[CH2]11C6H4SO3Na} (muối alkylbenzene sulfonate).
}
\end{ex}

\begin{ex}
Từ m kg mùn cưa chứa 60\% cellulose (còn lại là tạp chất trơ) sản xuất được 100 kg glucose với hiệu suất toàn bộ quá trình là 75\%. Giá trị của m là bao nhiêu?
\shortans{200}
\loigiai{
\textbf{Trả lời: 200}\\
$m_{\text{cellulose, lt}} = 100 \times \frac{162}{180} = 90\text{ kg}$.\\
Với $H = 75\% \Rightarrow m_{\text{cellulose, tt}} = \frac{90}{75\%} = 120\text{ kg}$.\\
Khối lượng mùn cưa: $m = \frac{120}{60\%} = \mathbf{200\text{ kg}}$.
}
\end{ex}

\vspace{6pt}
\noindent{\color{mainGreen}\textbf{\Large PHẦN IV. TỰ LUẬN}}
\vspace{4pt}

\noindent\textbf{Câu 1.} Hoàn thành dãy chuyển hóa sau (ghi rõ điều kiện nếu có):\\
\centerline{Tinh bột $\rightarrow$ Glucose $\rightarrow$ Ethanol $\rightarrow$ Acetic acid $\rightarrow$ Methyl acetate}

\vspace{6pt}
\multido{\i=1+1}{5}{\noindent\dotfill\\[4pt]}

\vspace{6pt}
\noindent\textbf{Câu 2.} Phản ứng tổng hợp glucose trong cây xanh cần được cung cấp năng lượng từ ánh sáng mặt trời:\\
\centerline{$6\ce{CO2} + 6\ce{H2O} + 673\text{ kcal} \xrightarrow[\text{diệp lục}]{\text{as}} \ce{C6H12O6} + 6\ce{O2}$}\\
Cứ trong một phút, mỗi $\text{cm}^2$ lá xanh nhận được $0{,}5\text{ cal}$ năng lượng mặt trời, nhưng chỉ có $10\%$ được sử dụng vào phản ứng tổng hợp glucose. Thời gian để một cây có 1000 lá xanh (diện tích mỗi lá $10\text{ cm}^2$) sản sinh được 18 gam glucose là bao lâu?

\vspace{6pt}
\multido{\i=1+1}{6}{\noindent\dotfill\\[4pt]}

\vspace{6pt}
\noindent\textbf{Câu 3.} Cho X là hexapeptide Ala-Gly-Ala-Val-Gly-Val và Y là tetrapeptide Gly-Ala-Gly-Glu. Thủy phân hoàn toàn m gam hỗn hợp gồm X và Y thu được 4 amino acid, trong đó có 30 gam glycine và 28,48 gam alanine. Giá trị của m là bao nhiêu?

\vspace{6pt}
\multido{\i=1+1}{7}{\noindent\dotfill\\[4pt]}

\begin{center}
\textbf{--- HẾT ---}
\end{center}
\end{document}
